\documentclass[12pt]{article}
\usepackage[utf8]{inputenc}
\usepackage[T1]{fontenc}
\usepackage{hyperref}
\title{Europanto Language Resources for K-12 Curriculum Development}
\author{Compiled by Claude Sonnet 5.5 (claude-sonnet-5-5) for Andr\'as Kornai}
\date{}
\begin{document}
\maketitle

\section{Introduction}
Europanto (ISO 639-3 eur) is a humorous macaronic jargon devised in 1996 by Diego Marani, a translator at the Council of the European Union. It has no fixed grammar and no native speakers. Marani wrote newspaper columns and a novel in it. This is a playful written practice, not a language that sustains a school curriculum, and the dossier is deliberately short.

\section{Primary Educational Resources}
No textbooks, primers or school programmes were found.

\section{Linguistic and Reference Materials}
\begin{itemize}
\item Wikipedia, Europanto: \url{https://en.wikipedia.org/wiki/Europanto}
\item World Wide Words on Europanto: \url{https://worldwidewords.org/tp-eur2.html}
\item Rosetta Panglossia entry: \url{https://rosettapanglossia.longnow.org/index.php?title=Eur}
\end{itemize}

\section{Religious and Community Resources}
None found.

\section{Digital and Audio Resources}
\begin{itemize}
\item Mediamatic page on Marani and Europanto: \url{https://www.mediamatic.net/en/page/8718/europanto}
\end{itemize}

\section{Curriculum Development Context}
At most Europanto could serve as a short enrichment unit on language contact, borrowing and multilingualism, using Marani's published columns. No evidence of existing school use was found.

\end{document}